\documentclass[12pt, a4paper]{article}
\usepackage[utf8]{inputenc}
\usepackage[T1]{fontenc}
\usepackage[french]{babel}
\usepackage{amsmath, amssymb, amsthm}
\usepackage{geometry}
\usepackage{listings}
\geometry{margin=2.5cm}

\newtheorem{theorem}{Théorème}[section]
\newtheorem{lemma}[theorem]{Lemme}
\newtheorem{definition}[theorem]{Définition}

\title{Analyse Structurale de la Conjecture d'Erdős-Straus}
\author{Charles EDOU NZE}
\date{\today}

\lstdefinelanguage{lean}{
  keywords={import, def, theorem, lemma, by, admit, Prop, Nat, open, section, Exists, fun},
  sensitive=true,
  comment=[l]--
}

\begin{document}
\maketitle

\begin{abstract}
Ce document présente une analyse rigoureuse de la conjecture d'Erdős-Straus, stipulant que pour tout $n \geq 2$, l'équation $4/n = 1/x + 1/y + 1/z$ admet des solutions entières strictement positives. Nous introduisons un système de congruences complètes, et définissons les solutions de Type A et B en s'appuyant sur les travaux de Lopez. \\ \vspace{1cm} \\ \textit{Charles EDOU NZE, chercheur indépendant}
\end{abstract}

\tableofcontents
\newpage

\section{Définitions Axiomatiques}
\begin{definition}
Soit $n \in \mathbb{N}$ avec $n \geq 2$. L'équation d'Erdős-Straus est l'équation diophantienne :
\[\frac{4}{n} = \frac{1}{x} + \frac{1}{y} + \frac{1}{z}\]
où $x, y, z \in \mathbb{Z}_{>0}$.
\end{definition}

\section{Recherche Littéraire et Stratégie}
La conjecture reste ouverte. Des travaux récents (Lopez) ont défini un système complet de congruences qui limite les exceptions potentielles, introduisant des solutions de Type A et de Type B selon leur structure modulaire. La preuve se décompose en lemmes isolant les classes de congruence modulo un entier donné.

\section{Démonstration et Lemmes}
\begin{lemma}
Pour tout $n \equiv 3 \pmod 4$, l'équation admet une solution entière positive.
\end{lemma}
\begin{proof}
Soit $n = 4k+3$ pour $k \in \mathbb{N}$. Posons $x = k+1$. Puisque $n \geq 3$, nous avons $k \geq 0$ et donc $x \geq 1$. Considérons la différence :
\begin{align*}
\frac{4}{n} - \frac{1}{x} &= \frac{4}{4k+3} - \frac{1}{k+1} \\
&= \frac{4(k+1) - (4k+3)}{(4k+3)(k+1)} \\
&= \frac{4k+4 - 4k - 3}{n(k+1)} \\
&= \frac{1}{n(k+1)}
\end{align*}
Nous utilisons l'identité égyptienne $1/A = 1/(A+1) + 1/(A(A+1))$ avec $A = n(k+1)$.
\begin{align*}
\frac{1}{n(k+1)} &= \frac{1}{n(k+1) + 1} + \frac{1}{n(k+1)(n(k+1) + 1)}
\end{align*}
En posant $y = n(k+1)+1$ et $z = n(k+1)(n(k+1)+1)$, nous obtenons $x,y,z \in \mathbb{Z}_{>0}$ satisfaisant l'équation.
\end{proof}

\section{Architecture pour Autoformalisation}
\begin{lstlisting}[basicstyle=\ttfamily\small, breaklines=true]
import Mathlib.Data.Nat.Basic
import Mathlib.Tactic.Omega
import Mathlib.Tactic.Ring

def SatisfiesErdosStraus (n : Nat) : Prop :=
  Exists (fun x => Exists (fun y => Exists (fun z => x > 0 /\ y > 0 /\ z > 0 /\ 4 * x * y * z = n * (y * z + x * z + x * y))))

lemma erdos_straus_mod_4_3 (k : Nat) : SatisfiesErdosStraus (4 * k + 3) := by
  admit

theorem erdos_straus_conjecture (n : Nat) (hn : n >= 2) : SatisfiesErdosStraus n := by
  admit
\end{lstlisting}

\end{document}
